\documentclass[conference]{IEEEtran}
\usepackage[utf8]{inputenc}
\usepackage[T1]{fontenc}
\usepackage{graphicx}
\usepackage{array}
\usepackage{booktabs}
\usepackage{amsmath}
\usepackage{url}
\usepackage{cite}

% Paper del trabajo final — Visión por Computador (Maestría en IA, Grupo 3)
% Formato: WVC main track — IEEE conference, máx. 6 páginas,
% título y abstract en inglés (abstract <= 150 palabras).
% Para revisión doble ciego de WVC real: comentar \author y usar el placeholder.

\begin{document}

% Logo on the LEFT of the title, single row (m{} column = vertical centering)
\title{\begin{tabular}{@{}m{2.3cm}@{\hspace{0.3cm}}m{\dimexpr\textwidth-3.0cm\relax}@{}}

\includegraphics[width=2.1cm]{figuras/logo_uni.png} &
\centering\arraybackslash Automatic Extraction of Electoral Results from the
Tally Sheets of the 2026 Peruvian General Elections
\end{tabular}}

\author{
\IEEEauthorblockN{Josemanuel Rossy Ca\~nari Palante, Kenny Asto Hinostroza,
Melissa Dessire Aylas Barranca, Carlos P\'erez P\'erez}
\IEEEauthorblockA{Master's Program in Artificial Intelligence --- Computer Vision, Group 3\\
Universidad Nacional de Ingenier\'ia, Lima, Peru}
}
% \author{\IEEEauthorblockN{Anonymous WVC submission}} % versión doble ciego

\maketitle

\begin{abstract}
Official electoral results in Peru are published as scanned tally sheets
(\emph{actas}), which prevents large-scale automatic processing. We present an
end-to-end pipeline that downloads actas from the official ONPE portal,
locates 46 fields per document through a fixed template anchored by
fiducial-mark registration, and reads handwritten vote counts with
digit-restricted OCR. Evaluation is annotation-free: ground truth comes from
the officially digitized results published by the same portal, and the test
set is a seeded random national sample (100 actas, 23 departments). The
population is heterogeneous: 24\% are born-digital documents whose text layer
yields 95.1\% field accuracy without vision. On 76 scanned
handwritten actas (3,268 fields), the pipeline attains 59.6\% field accuracy
(95\% CI [56.6, 62.5]) and 0.401 CER, against 48.4\% for an end-to-end
multimodal LLM prompted few-shot, quantifying the ceiling of generic
recognition and motivating fine-tuned recognizers trained on the ${\sim}43$
field--value pairs each acta provides for free.
\end{abstract}

\begin{IEEEkeywords}
document analysis, optical character recognition, handwritten digit
recognition, electoral documents, template registration
\end{IEEEkeywords}

\noindent\textbf{\textit{Resumen}}---Los resultados electorales oficiales del
Per\'u se publican como actas de escrutinio escaneadas, lo que impide su
procesamiento autom\'atico a gran escala. Presentamos un pipeline de extremo a
extremo que descarga las actas del portal oficial de la ONPE, localiza 46
campos por documento mediante una plantilla fija anclada por registro de
marcas fiduciales, y lee los conteos manuscritos con OCR restringido a
d\'igitos. Dos decisiones de dise\~no hacen la evaluaci\'on reproducible y sin
anotaci\'on manual: la verdad de terreno proviene de los resultados
oficialmente digitados que publica el mismo portal, y el conjunto de prueba es
una muestra aleatoria nacional con semilla publicada (100 actas, 23
departamentos). La poblaci\'on de actas es heterog\'enea: el 24\% son
documentos nacidos digitales cuya capa de texto rinde 95.1\% de exactitud por
campo sin visi\'on alguna. Sobre las 76 actas manuscritas escaneadas (3,268
campos), el pipeline alcanza 59.6\% de exactitud por campo (IC 95\% [56.6,
62.5]) y CER de 0.401, frente a 48.4\% de un LLM multimodal de extremo a
extremo con few-shot, cuantificando el techo del reconocimiento gen\'erico y
motivando reconocedores afinados con los ${\sim}43$ pares campo--valor que
cada acta provee gratuitamente.\\[4pt]
\noindent\textbf{\textit{Palabras clave}}---an\'alisis de documentos,
reconocimiento \'optico de caracteres, reconocimiento de d\'igitos
manuscritos, documentos electorales, registro de plantilla.

\section{Introduction}
Scrutiny tally sheets (\emph{actas de escrutinio}) are the official source of
electoral results for each polling station in Peru. Although the National
Office of Electoral Processes (ONPE) publishes every acta digitally, the
information they carry is visual: a scanned form where members of the polling
station wrote vote counts by hand. Manual extraction is slow, error-prone and
does not scale to the ${\sim}88{,}000$ polling stations of a general election,
which limits automatic auditing and independent verification.

This work asks a concrete question: \emph{how far can a classical computer
vision pipeline go on real, nationally sampled Peruvian actas, and what does
that imply for the design of a production system?} We build and evaluate a
five-stage pipeline (acquisition $\rightarrow$ preprocessing $\rightarrow$
region detection $\rightarrow$ recognition $\rightarrow$ structured output)
on the first round of the 2026 General Elections. Our contributions are:

\begin{enumerate}
\item \textbf{An annotation-free, reproducible evaluation protocol.} The same
portal that publishes each acta image also publishes its officially digitized
vote counts through an internal API. We use those values as ground truth,
which turns every downloaded acta into ${\sim}43$ labeled field--value pairs
at zero annotation cost, and we sample the test set randomly at national
scale with a published seed.
\item \textbf{Template detection made robust by fiducial registration.} The
acta is a fixed-layout form with printed registration marks. We show that
fixed fractional templates break under real-world scanner framing (black
bands, offsets) and that a RANSAC-based translation matching of fiducial
marks followed by a partial affine warp recovers alignment, adding $+8.7$
accuracy points.
\item \textbf{A population finding with architectural consequences.} In a
random national sample, 24\% of actas are born-digital (STAE) documents whose
PDF text layer can be parsed directly with 95.1\% field accuracy; only 76\%
are scanned handwritten forms that require vision. A production system should
classify the acta type first and route accordingly.
\item \textbf{A quantified ceiling for generic OCR.} On handwritten actas,
generic digit-restricted OCR with careful cell cleaning reaches 59.6\% field
accuracy (95\% CI [56.6, 62.5], clustered bootstrap). This bounds what can be
expected without training and motivates fine-tuning transformer recognizers
with the free labels described above.
\item \textbf{A head-to-head comparison with an end-to-end multimodal LLM.}
Under an identical protocol, few-shot prompting of a current multimodal LLM
reaches 48.4\% against the pipeline's 60.9\% on the same actas, which places
its value in zero-annotation bootstrapping rather than in recognition
accuracy.
\end{enumerate}

% Para un envío doble ciego a WVC, comentar la nota al pie siguiente.
Code, per-field evaluation artifacts and the sampling seed are
public.\footnote{\url{https://github.com/carlosperez100/onpe_actas}}

\section{Related Work}
\textbf{Real-time detection.} RT-DETR \cite{rtdetr} showed that a carefully
designed DETR outperforms comparable-scale YOLO detectors (L/X) in both
accuracy and speed while
removing non-maximum suppression (NMS) and its manually tuned thresholds. For
tally sheets --- tables of contiguous cells frequently overlapped by stamps
and signatures --- threshold-free, end-to-end detection is attractive, and
RT-DETR's end-to-end latency protocol is the one we will adopt to compare
learned detectors (YOLOv11 vs.\ RT-DETR) in future work. In this paper, detection is
solved by template registration, which doubles as an automatic label
generator for those detectors.

\textbf{Text recognition.} TrOCR \cite{trocr} established transformer
encoder--decoder OCR with pre-trained language models, and DTrOCR
\cite{dtrocr} showed that a decoder-only generative architecture surpasses it
on printed, scene and handwritten text (CER 2.38 on IAM without an external
language model). These results justify transformer recognizers as the upgrade
path for our handwritten-digit fields; a caveat for our setting is that
language priors may ``correct'' digit strings toward plausible sequences,
which matters when reading isolated numbers.

\textbf{Unified document parsing.} OmniParser \cite{omniparser} unifies text
spotting, key information extraction and table recognition through
structured-point sequences, and highlights the weakness of monolithic
sequence generation: Donut \cite{donut} collapses on long tables due to
error accumulation, dropping to 22.7 TEDS on PubTabNet vs.\ 88.83 for
OmniParser. An acta is precisely a long table plus key
fields; OmniParser's explicit localization of every extracted value also
matches the auditability requirement of electoral documents. We keep a
modular pipeline --- whose per-stage metrics are easier to diagnose --- and
treat unified end-to-end parsing as the alternative to compare against.

None of these works evaluates handwritten digits in Spanish on official forms
with stamps and signatures; that combination is the empirical gap this
project addresses.

\section{Dataset}
\subsection{Source and acquisition}
All documents come from the official ONPE results portal for the 2026
General Elections (first round, April 12, 2026). The portal exposes an
internal API: searching for a polling station returns its actas and their processing statuses;
the acta detail lists its files, where type 1 is the scrutiny sheet; a
signed S3 URL serves the PDF. Downloads use browser TLS impersonation, a
conservative request rate, and a resumable manifest
(\texttt{src/scraper/} in our repository).

\subsection{Sampling design}
The universe holds ${\sim}88{,}064$ station codes (upper bound located by
exponential-plus-binary search against the API). We drew a simple random
sample without replacement with a \textbf{published seed (2026)}: 100 actas,
100\% hit rate, covering 23 of Peru's department-level \emph{ubigeo} prefixes (the Peruvian standard geographic location code) with a
size-proportional distribution (Lima concentrates 36\%, consistent with its
electorate share). This corrects the pilot bias of sequential station codes
from a single district.

\subsection{Population heterogeneity}
Classifying each PDF by page count and text layer reveals two document types:
\textbf{76 handwritten actas} (single-page scans, empty text layer) and
\textbf{24 born-digital STAE actas} (two-page digitally signed documents with
typeset vote counts and an extractable text layer). The mix itself is a
sample estimate: born-digital prevalence is 24\% with a 95\% Wilson interval
of [16.7, 33.2] --- enough to establish that the population is mixed, not
enough to pin the exact share. All vision experiments
run on the 76 handwritten actas; STAE actas are processed by direct text
parsing (Sec.~\ref{sec:stae}).

\subsection{Ground truth}
For every sampled station we store the officially digitized counts returned
by the API: votes per political organization (38 rows), blank/null/challenged
votes, totals and electorate. Form positions 04 and 14 correspond to
withdrawn parties, are empty on the form and null in the API; they are
excluded (a null-to-null match would spuriously inflate accuracy). This
yields \textbf{43 evaluable fields per acta, 3,268 in total}.

\section{Method}
\subsection{Pipeline overview}
(1) \emph{Acquisition}: PDF download plus official ground truth.
(2) \emph{Preprocessing}: 300 DPI rasterization, deskew by dominant angle,
CLAHE contrast normalization, non-local-means denoising (OpenCV).
(3) \emph{Region detection}: template with fiducial registration (below).
(4) \emph{Recognition}: digit-restricted OCR on cleaned cells.
(5) \emph{Structured output}: one JSON per acta with the proposal's schema.

\subsection{Template regions with fiducial registration}
\label{sec:registro}
The acta is a fixed-layout form, so 46 regions (38 vote cells arranged as
equally spaced rows, four summary rows, station number, electorate, citizen
total, observations) are defined once as page fractions calibrated on a
reference acta (Fig.~\ref{fig:overlay}). National sampling exposed the
failure mode of fixed fractions: scanner framing varies (black margin bands,
content offsets of up to ${\sim}184$ px), shifting every cell one row off and
producing CER $>$ 1. The form itself provides the fix: printed black
registration squares along the border. We detect them as dark, square,
high-fill connected components inside border bands; match detected marks to
15 calibrated reference positions by translation RANSAC (each detected--%
reference pair proposes a shift, the shift with most inliers wins), robust to
missing marks and to false positives such as the printed sample digits; and
estimate a partial affine transform that warps the page to the reference
frame before cropping.

\begin{figure}[t]
\centering
\includegraphics[width=\columnwidth]{figuras/plantilla_overlay.png}
\caption{Reference acta with the 46 template regions: 38 vote cells (blue),
summary and header fields (red). Region coordinates are page fractions,
applied after fiducial registration.}
\label{fig:overlay}
\end{figure}

\subsection{Recognition of handwritten counts}
Cells are read with EasyOCR restricted to a digit vocabulary, after a
cleaning cascade motivated by error analysis on a 10-acta pilot
(Fig.~\ref{fig:limpieza}): (i) an interior margin removes printed cell
borders, which were previously misread as extra digits; (ii) per-cell Otsu binarization isolates
pen ink and survives the alternating shaded rows of the form; (iii) the
dotted intra-cell separators --- as dark as ink, hence immune to
thresholding --- are erased \emph{by template position} (fixed fractions of
the cell width); (iv) a connected-component speckle filter removes
digitization noise; (v) fields with per-digit printed boxes (citizen total)
are read box by box. A fallback re-reads the raw crop when cleaning empties
faint strokes.

\begin{figure}[t]
\centering
\includegraphics[width=\columnwidth]{figuras/limpieza_celdas.png}
\caption{Ink isolation on problem cells (top: crop; bottom: cleaned). Dotted
separators and printed borders --- previously read as digits (e.g.\
$18\!\rightarrow\!418$) --- are removed by threshold, template-positioned
erasure and speckle filtering.}
\label{fig:limpieza}
\end{figure}

\subsection{Domain rule: empty cell means zero}
Polling-station members frequently leave the cell blank when a party receives
no votes. A blank cell is not an OCR failure but form semantics, so a
post-processing rule maps unread party cells (and challenged votes) to zero.
We report it separately to avoid crediting vision with domain knowledge.

\subsection{Routing born-digital actas}
\label{sec:stae}
STAE actas are parsed directly from the PDF text layer: each party name is
followed by its typeset count, and summary labels precede their values.
Matching extracted names to official ones by normalized string equality
yields field accuracy without any image processing.

\section{Experiments}
\subsection{Metrics and statistical protocol}
\textbf{Field accuracy}: exact match between the read integer and the
official value. \textbf{CER}: character error rate between the digit strings.
Fields within an acta share writer and scan quality, so the acta is the
statistical cluster: confidence intervals use a \textbf{clustered bootstrap}
(10,000 resamples of whole actas). All evaluation artifacts (per-field
errors, per-acta metrics) ship with the repository.

\subsection{Pilot ablation (10 actas, 430 fields)}
The pilot --- sequential stations of one district, hence biased but useful
for component analysis --- quantifies each recognition improvement
(Table~\ref{tab:pilot}).

\begin{table}[t]
\caption{Pilot ablation on 10 actas (430 fields), field accuracy and CER.}
\label{tab:pilot}
\centering
\begin{tabular}{lrr}
\toprule
Configuration & Acc. (\%) & CER \\
\midrule
v1: raw digit OCR & 43.95 & 0.524 \\
v2: + ink isolation, borders, sub-boxes & 52.09 & 0.419 \\
v3: + separator erasure, Otsu, speckle & 52.79 & 0.476 \\
v4: + empty-cell-is-zero rule & 58.37 & 0.420 \\
\bottomrule
\end{tabular}
\end{table}

\subsection{National evaluation (76 handwritten actas, 3,268 fields)}
Table~\ref{tab:national} shows the same system on the random national sample.
Fiducial registration (Sec.~\ref{sec:registro}) is decisive: without it,
framing variation severely degrades part of the sample (worst acta 14\%,
CER $>$ 1). Fiducial registration adds $+8.7$ points and lifts the
worst acta from 14\% to 30\%; the domain rule adds $+21.6$ points ---
nationally, blank zero-cells are far more common than in the pilot district.

\begin{table}[t]
\caption{National random sample, 76 handwritten actas (3,268 fields).}
\label{tab:national}
\centering
\begin{tabular}{lrr}
\toprule
Configuration & Acc. (\%) & CER \\
\midrule
Fixed template, raw OCR & 26.81 & 0.724 \\
Fixed template + zero rule & 50.86 & 0.482 \\
Registered (v5), raw OCR & 38.00 & 0.617 \\
\textbf{Registered (v5) + zero rule} & \textbf{59.58} & \textbf{0.401} \\
\bottomrule
\end{tabular}
\end{table}

\textbf{Main result.} Field accuracy \textbf{59.58\%, 95\% CI [56.61, 62.48]}
($\pm 2.94$ points); CER \textbf{0.401, 95\% CI [0.360, 0.444]}. Per-acta
accuracy averages 59.6\% (SD 13.0, range 30--84\%): residual variance stems
from handwriting and digitization quality, no longer from alignment. The
pilot figure (58.4\%) happens to be close, but only the national estimate
has a valid margin of error and offers national coverage.

\textbf{Born-digital actas.} Text-layer parsing of the 24 STAE actas reaches
\textbf{95.12\%} (955/1,004 fields); the remaining errors stem from
name matching, not from reading. Twenty-four actas suffice here because
parsing is deterministic: between-acta variance is below binomial
(ICC $\approx 0$, design effect 1.0), so the clustered bootstrap interval is
[94.5, 95.7], nine times tighter than the handwritten estimate, whose
design effect is 3.07 for a comparable field count. The 35.5-point gap
(95\% CI [32.6, 38.5]) quantifies the value of routing by document type.

\subsection{End-to-end multimodal LLM baseline}
\label{sec:llm}
A natural alternative to the modular pipeline is to hand the whole page to a
multimodal LLM together with a few solved examples, letting it localize and
read in a single step. We evaluate that baseline\footnote{Gemini 2.5 Flash,
free tier, accessed July 2026. Code: \texttt{src/recognition/llm\_ocr.py}.}:
the model receives the full acta image plus two image--JSON exemplars and must
return the vote-count JSON, with no template, no registration and no OCR
engine. Both systems are scored on the same 15 handwritten actas (645 fields)
with the same protocol (Table~\ref{tab:llm}).

\begin{table}[t]
\caption{Modular pipeline vs.\ end-to-end LLM few-shot on the same 15
handwritten actas (645 fields).}
\label{tab:llm}
\centering
\begin{tabular}{lrr}
\toprule
Method & Acc. (\%) & CER \\
\midrule
\textbf{Pipeline (registered template + digit OCR)} & \textbf{60.93} & \textbf{0.395} \\
LLM few-shot, end-to-end (2 exemplars) & 48.37 & 0.591 \\
\bottomrule
\end{tabular}
\end{table}

The pipeline wins by 12.6 points. The LLM parses the form structure
competently, but it is not trained for handwritten Peruvian digits, and it
degrades faster on poor scans: on the two worst actas it drops to 20.9\% and
23.3\% where the pipeline still reads 34.9\% and 53.5\%, since cropping gives
the recognizer an isolated cell while the LLM must fight the surrounding
noise. This is consistent with DTrOCR \cite{dtrocr}: strong handwriting
recognition comes from task-specific training, not from general capability.
Its advantage is orthogonal to accuracy --- it required zero detection
engineering, whereas the pipeline needed template calibration and fiducial
registration --- which makes it useful as a bootstrap annotator for a new
form layout, and as a pre-labeler whose output a human corrects, rather than
as a replacement for a specialized recognizer.

\subsection{Error analysis}
Remaining errors concentrate in: isolated thin digits missed by the detector
stage of the OCR engine; strokes invading the erased separator bands;
shape confusions ($8{\leftrightarrow}18$, $7{\leftrightarrow}9$) on wide
handwriting; and multi-box fields. These are recognition errors, not
localization errors --- exactly the component that fine-tuning replaces.

\section{Conclusions and Future Work}
A classical pipeline over real, nationally sampled Peruvian actas achieves
59.6\% field accuracy (95\% CI [56.6, 62.5]) on handwritten documents and
95.1\% on born-digital ones, with an evaluation protocol that needs no manual
annotation. Three design lessons stand out: register templates on the form's
fiducial marks before any fixed-coordinate cropping; separate domain
semantics (blank $=$ 0) from vision performance; and classify the document
type before choosing a processing path. A fourth lesson comes from the LLM
baseline: prompting a general multimodal model end-to-end is cheaper to build
but 12.6 points less accurate (48.4\%), so on this task generality does not
yet substitute for specialization. The quantified ceiling of generic
OCR motivates the next steps: fine-tune a transformer recognizer
\cite{trocr,dtrocr} on the free field--value pairs (${\sim}43$ per acta,
${\sim}43{,}000$ from 1,000 actas), train learned detectors (YOLOv11,
RT-DETR \cite{rtdetr}) on labels auto-generated by the registered template,
compare the modular pipeline against unified parsers \cite{omniparser,donut},
and use the LLM as a pre-labeler in the human-correction loop that produces
the fine-tuning set.

\begin{thebibliography}{9}
\bibitem{rtdetr} Y.~Zhao, W.~Lv, S.~Xu, J.~Wei, G.~Wang, Q.~Dang, Y.~Liu,
and J.~Chen, ``DETRs beat YOLOs on real-time object detection,'' in
\emph{Proc. IEEE/CVF CVPR}, 2024, pp. 16965--16974.
\bibitem{dtrocr} M.~Fujitake, ``DTrOCR: Decoder-only transformer for optical
character recognition,'' in \emph{Proc. IEEE/CVF WACV}, 2024, pp. 8025--8035.
\bibitem{omniparser} J.~Wan, S.~Song, W.~Yu, Y.~Liu, W.~Cheng, F.~Huang,
X.~Bai, C.~Yao, and Z.~Yang, ``OmniParser: A unified framework for text
spotting, key information extraction and table recognition,'' in \emph{Proc.
IEEE/CVF CVPR}, 2024, pp. 15641--15653.
\bibitem{trocr} M.~Li, T.~Lv, J.~Chen, L.~Cui, Y.~Lu, D.~Florencio, C.~Zhang,
Z.~Li, and F.~Wei, ``TrOCR: Transformer-based optical character recognition
with pre-trained models,'' in \emph{Proc. AAAI}, 2023, pp. 13094--13102.
\bibitem{donut} G.~Kim, T.~Hong, M.~Yim, J.~Nam, J.~Park, J.~Yim, W.~Hwang,
S.~Yun, D.~Han, and S.~Park, ``OCR-free document understanding transformer,''
in \emph{Proc. ECCV}, 2022, pp. 498--517.
\end{thebibliography}

\end{document}
